\subsection{From Positional Access to Hash-Based Access}

Sequential types (\texttt{str}, \texttt{list}, \texttt{tuple}) expose \texttt{\_\_getitem\_\_}: an integer index selects a slot in an ordered component array. Hash-based types (\texttt{set}, \texttt{dict}) reject positional indexing because placement is driven by hash coordinates, not ordinal rank.

\begin{quote}
Sequential containers locate components by position. Hash-based containers locate components by hash-derived table slots, confirmed with equality.
\end{quote}

\subsubsection{Hash-Based Collections: Why Sets and Dictionaries Locate Components by Hash-Derived Table Slots}

A \texttt{set} answers membership; a \texttt{dict} maps a hashable key to an associated value. Neither type implements meaningful integer \texttt{\_\_getitem\_\_} access: \texttt{set} lacks subscripting entirely; \texttt{dict} subscripting takes a \textbf{key}, not an index.

Lookup follows the same C-level route in both cases:

\begin{verbatim}
set:  element -> hash(element) -> probe table -> __eq__ confirmation
dict: key     -> hash(key)     -> probe table -> __eq__ confirmation -> value
\end{verbatim}

The hash integer is a search coordinate, not the stored object. Membership and key retrieval therefore avoid $O(n)$ sequential scans when the table load factor stays within design bounds.
